\documentclass[12p, oneside]{book}
\usepackage[T1]{fontenc}
\usepackage[polish]{babel}
\usepackage{color}
\usepackage[table, usenames, dvipsnames]{xcolor}
\usepackage{hyperref}
\usepackage{graphicx}
\hypersetup{
    colorlinks=true, 
    linktoc=all,    
    linkcolor=blue,
}



\newcommand{\steno}[1]{\textbf{\textsf{\textcolor{violet}{ #1}}}}

\newcommand{\stenoExc}[1]{\textbf{\textsf{\textcolor{red}{ #1}}}}

\newenvironment{tabela2}
	{
		\rowcolors{1}{cyan}{Emerald}
		\begin{tabular}{c c}
	}
	{
		\end{tabular}
	}


\setcounter{secnumdepth}{0}
\begin{document}

\frontmatter

\title{Polski System Stenotypii}

\author{B. P.}
\maketitle
\newpage

\tableofcontents

\mainmatter

\chapter{Wstęp}

\chapter{Kurs}

\chapter{Pełny opis systemu}
\section{Klawiatura}
\includegraphics[width=10cm]{grafika/pss/klawiatura.jpg}

Powyżej przedstawiony jest układ klawiatury wykorzystywany przez PSS. Klawiatura jest podzielona na cztery części:
\begin{enumerate}
	\item Lewa górna --- odpowiedzialna za spółgłoski początkowe oraz samogłoski posiłkowe --- tę część obsługuje się palcami lewej dłoni
	\item Lewa dolna --- klawisz J odpowiedzialny za udźwięcznianie oraz dwa klawisze samogłosek --- tę część obsługuje się lewym kciukiem
	\item Prawa dolna --- trzy klawisze samogłosek --- tę część obsługuje się prawym kciukiem
	\item Prawa górna --- odpowiedzialna za spółgłoski końcowe oraz samogłoski dodatkowe --- tę część obsługuje się palcami prawej dłoni
\end{enumerate}


\section{Budowa akordu}
\subsection{Rdzenie}
\subsection{Przedrostki}
\subsection{Przyrostki}
\subsection{Specjalne}
\section{Samogłoski}
\section{Spółgłoski w nagłosie}
\section{Spółgłoski w wygłosie}

\end{document}
